\SourceLeaf{242}%
 \(x,y,z\) als Funktionen der Unbestimmten \(p,q\) darstellen. Der allgemeine Ausdruck für die Größe eines beliebigen Linienelements genügt. Gehen wir nun zu einigen Anwendungen dieses überaus wichtigen Satzes über.

Nehmen wir an, unsere gekrümmte Fläche lasse sich auf eine andere, gekrümmte oder ebene Fläche abwickeln, sodass jedem durch die Koordinaten \(x,y,z\) bestimmten Punkt der ersten Fläche ein bestimmter Punkt der zweiten Fläche mit den Koordinaten \(x',y',z'\) entspricht. Offenbar lassen sich also auch \(x',y',z'\) als Funktionen der Unbestimmten \(p,q\) auffassen. Für das Element \(\BareRoot{(\mathrm{d}x'^2+\mathrm{d}y'^2+\mathrm{d}z'^2)}\) ergibt sich daher ein Ausdruck der Form
\[
\BareRoot{(E'\mathrm{d}p^2+2F'\mathrm{d}p\mathord{.}\mathrm{d}q+G'\mathrm{d}q^2)}
\]
wobei auch \(E',F',G'\) Funktionen von \(p,q\) bezeichnen. Schon aus dem Begriff der \textit{Abwicklung} einer Fläche auf eine andere ergibt sich, dass entsprechende Elemente auf beiden Flächen notwendig gleich sind. Daraus ergeben sich die Identitäten
\[
E=E',\quad F=F',\quad G=G'
\]
Die Formel des vorigen Artikels führt somit unmittelbar zu dem bemerkenswerten

\textsc{Satz.} \textit{Wird eine gekrümmte Fläche auf irgendeine andere Fläche abgewickelt, so bleibt das Krümmungsmaß in jedem einzelnen Punkt unverändert.}

Offenbar behält auch \textit{jeder endliche Teil der gekrümmten Fläche nach der Abwicklung auf eine andere Fläche dieselbe Integralkrümmung.}

Die auf eine Ebene abwickelbaren Flächen bilden den Spezialfall, auf den die Geometer ihre Untersuchungen bisher beschränkt haben. Unsere Theorie zeigt unmittelbar, dass das Krümmungsmaß solcher Flächen in jedem Punkt \(=0\) ist. Wird ihre Beschaffenheit daher nach der dritten Methode beschrieben, so gilt überall
\[
\frac{\mathrm{d}\mathrm{d}z}{\mathrm{d}x^2}\mathord{.}\frac{\mathrm{d}\mathrm{d}z}{\mathrm{d}y^2}-(\frac{\mathrm{d}\mathrm{d}z}{\mathrm{d}x\mathord{.}\mathrm{d}y})^2=0
\]
Dieses Kriterium ist zwar seit Langem bekannt, wird aber zumeist, zumindest nach unserem Urteil, mit geringerer Strenge bewiesen, als man sich wünschen könnte.
\Article{13}
Was wir im vorigen Artikel ausgeführt haben, hängt mit einer besonderen Betrachtungsweise von Flächen zusammen, die es in hohem Maße verdient, von Geometern sorgfältig weiterentwickelt zu werden. Betrachtet man nämlich eine Fläche losgelöst von ihrer Rolle als Grenze eines Körpers und fasst sie als einen
